% ---------------------------------------------------------------------
% STATUT : RÉDIGÉ (transplanté de exploratory/cascade_qualitative/
%          chapitre_robustesse.tex, préambule retiré).
% RESTE À FAIRE : ajouter une ligne au tableau des paramètres pour u_ij
%          (statut : posé / élicité, sensibilité mesurée par l'échelle
%          de repli du chapitre 5).
% ---------------------------------------------------------------------
\chapter{Inventaire des hypothèses et de leur sensibilité}
\label{ch:robustesse}

\noindent Le chapitre~\ref{chap:resultats} a conduit les tests de robustesse eux-mêmes
(tornado, sensibilités des hypothèses restantes, invariance de l'ordre de remédiation). Ce
chapitre ne les refait pas : il en donne la \textbf{synthèse en une table}, puis l'inventaire
exhaustif des paramètres, de leur nature et de leur sensibilité. Il sert de pièce
justificative, de sorte qu'aucun chiffre du mémoire ne soit sans origine déclarée.

\begin{cle}
\textbf{Thèse.} Les conclusions du modèle ne reposent pas sur les paramètres non calibrés.
Je le montre en deux temps : (i)~elles sont \textbf{insensibles} aux choix de modélisation
innocents (structure de la cascade, famille de dépendance, couche temporelle, source de
sévérité) ; (ii)~elles ne \textbf{réagissent} qu'aux leviers qui doivent compter (conformité,
contagion), ce qui est une \emph{qualité}, pas un défaut, c'est la sensibilité au risque que
la Formule Standard n'a pas. Reste une incertitude irréductible, la queue de sévérité, que
j'assume en présentant le SCR comme une \textbf{distribution}, jamais un point.
\end{cle}

\section{L'argument de robustesse, en une table}

\begin{center}\footnotesize
\renewcommand{\arraystretch}{1.3}
\begin{tabular}{@{}p{4.5cm} p{4.6cm} p{4.4cm} c@{}}
\toprule
\textbf{Conclusion} & \textbf{Ce qu'on fait varier} & \textbf{Effet mesuré} & \textbf{Verdict} \\
\midrule
La non-conformité coûte cher ($\Delta_{\text{DORA}}>0$, grand) &
source de sévérité (OpRisk / PRC), famille de copule ($\theta\in[1,4]$) &
facteur ${\sim}1{,}9$ entre sources ; $\Delta_{\text{DORA}}$ varie $<1{,}4\,\%$ sur $\theta$ &
\rob \\
Priorité de remédiation P1 puis P4 (les \emph{sources}) &
cadre de dépendance (cascade vs copule), branchement vs marche &
P1 et P4 en tête dans tous les cadres ; SCR $\pm 4\,\%$ (branchement) &
\rob \\
C'est l'\textbf{asymétrie} de $W$ qui porte le résultat &
knockout de l'asymétrie (3\,000 tirages), placebo directionnel &
sans la direction, le signal tombe sous le bruit du placebo &
\rob \\
Pas de couche de Hawkes &
noyau exponentiel puis à retard, logique two-phase &
excitation intra-journalière ; endogénéité $\to 0$ si exogène isolé &
\rob \\
Le SCR \emph{réagit} à la conformité et à la contagion &
score de conformité $q\in[0,1]$, gain $g\in[0{,}05\,;0{,}90]$ &
$-53\,\%$ de NC à C ; $\times 1{,}6$ sur la plage de $g$ &
\feat \\
Le SCR est une distribution large, pas un point &
incertitude de paramètre (bootstrap deux niveaux) &
facteur ${\sim}2{,}5$ sur la VaR par la seule incertitude de calibration &
\ass \\
\bottomrule
\end{tabular}
\end{center}

\vspace{4pt}
\textbf{Lecture.} Les quatre premières lignes sont l'\textbf{ossature} : le signe et l'ordre
de grandeur de l'effet DORA, la hiérarchie des piliers, le rôle de la direction et le rejet
du Hawkes tiennent quand on change ce qui pourrait sembler arbitraire. La cinquième ligne
n'est pas une fragilité mais le \textbf{c\oe ur du modèle} : un SCR qui ne bougerait pas avec la
conformité serait inutile (c'est le reproche fait à la Formule Standard,
chapitre~\ref{chap:resultats}). La dernière est la limite honnête : la queue de sévérité
domine l'incertitude, donc toute lecture ponctuelle du SCR serait une surinterprétation.

\section{Tableau des paramètres non calibrés (et de leur sensibilité)}

Trois natures sont distinguées : \textbf{calibré} sur données (avec son incertitude),
\textbf{posé} par jugement d'expert ou choix normatif (non identifiable, cf.\
chapitre~\ref{ch:identifiabilite}), et \textbf{semi-ancré} sur un proxy externe non
spécifique au périmètre.

\vspace{0.25cm}
\begin{center}\footnotesize
\renewcommand{\arraystretch}{1.28}
\begin{tabular}{@{}p{2.7cm} p{2.6cm} p{1.7cm} p{4.0cm} p{3.6cm}@{}}
\toprule
\textbf{Paramètre} & \textbf{Valeur} & \textbf{Nature} & \textbf{Justification} & \textbf{Sensibilité du SCR} \\
\midrule
Queue de sévérité $\xi$ & $0{,}90$ (OpRisk $0{,}60$, PRC $1{,}03$) & calibré, incertain & GPD sur pertes réelles ; hétérogène par catégorie ($0{,}92$--$1{,}37$) & \textbf{forte} : $0{,}68$--$0{,}93$ selon le seuil ; pilote la queue \\
Échelle GPD $\sigma$, seuil $u$, $p_u$ & $57{,}97$ M\euro{} ; $20{,}03$ M\euro{} ; $0{,}15$ & calibré & POT au percentile 85, conversion en euros & modérée (via la queue) \\
Fréquence $\lambda$ & ${\sim}21{,}6$/an (OpRisk) & calibré & incidents TIC matériels / 27 ans & linéaire sur le niveau \\
Sur-dispersion $\phi$ & $9{,}20$ & semi-calibré & Var/Moyenne observé, maintenu entre scénarios & faible \\
Choc commun fréquence $a$ & $0{,}60$ & \textbf{posé} (prudent) & dépasse l'IC empirique haut ($0{,}22$) : conservateur & enveloppe le modèle à paquets ($\times 2{,}5$) \\
Structure $W = g\,\texttt{TRANS}/\max_k s_k$ & classeur qualitatif & \textbf{posé} / normatif & jugement d'expert, à \textbf{éliciter} (Cooke) ; direction non identifiable & l'asymétrie porte le résultat (knockout) \\
Gain de propagation $g$ & $0{,}90$ base ; par état $0{,}45/0{,}68/0{,}90$ & \textbf{posé} & amplitude de $W$, bornée par Leontief $\rho(W)\le g<1$ & \textbf{levier} : $\times 1{,}6$ sur $[0{,}05\,;0{,}90]$ \\
Score de conformité $q$ / états & continu $[0,1]$ (ou NC/PC/C) & \textbf{posé} & décrit le profil, subsume 3 ou 5 niveaux & \textbf{levier} : $-53\,\%$ de NC à C \\
Multiplicateur de matérialité par état & $0{,}85/1{,}00/1{,}20$ & \textbf{posé} & module $p_u$ selon la conformité & modérée (canal détection) \\
Charge systémique $\gamma$ & $0{,}68$ & semi-ancré & proxy concentration cloud (AWS+Azure+GCP, Synergy 2025) & modérée \\
Seuil de crise $\theta_{\text{crise}}$ & $-2{,}5$ & \textbf{posé} & déclenche le régime de co-mouvement extrême & agit sur la seule queue crise \\
Probas d'état $P(\text{NC/PC/C})$ & $0{,}35/0{,}35/0{,}30$ & semi-ancré & enquêtes conformité DORA ${\sim}50\,\%$ non conformes (Deloitte, ESA) & déplace le mélange, pas les états \\
Copule de dépendance $\theta$ & Gumbel $1{,}8$ & \textbf{posé} & queue supérieure entre briques & \textbf{négligeable} : $\Delta_{\text{DORA}}$ $<1{,}4\,\%$ sur $[1,4]$ \\
Plafond de sévérité & $40$ M\euro{} & \textbf{posé} & capacité de réassurance ($\xi\ge 1$) & borne la queue extrême \\
\bottomrule
\end{tabular}
\end{center}

\vspace{4pt}
\textbf{Ce que la table rend visible.} Les seuls paramètres à \textbf{forte} sensibilité sont
soit calibrés sur données (la queue $\xi$, dont on assume et chiffre l'incertitude), soit des
\textbf{leviers voulus} ($g$, $q$) qui traduisent précisément l'objet du mémoire (l'effet de
la conformité et de la contagion sur le capital). Les paramètres purement \textbf{posés} qui
ne sont pas des leviers ($\theta$ de copule, $\gamma$, $\theta_{\text{crise}}$, probas d'état,
matérialité) ont une influence \textbf{faible à modérée} et documentée. Aucune conclusion ne
repose sur un paramètre à la fois non calibré \emph{et} déterminant hors de son rôle attendu.

\section{Les trois étages de l'incertitude du SCR}
\label{sec:trois-etages}

L'incertitude sur le SCR ne se réduit pas au paramètre. Elle se cumule sur trois étages, qu'il
faut nommer plutôt que confondre (figure~\ref{fig:bande-modele}). Le premier est l'incertitude
de \textbf{paramètre} : l'intervalle de confiance de $\xi$ à lui seul dilate la VaR d'un facteur
$2{,}6$ (chapitre~\ref{chap:socle}). Le deuxième est l'incertitude d'\textbf{identification} :
la direction de $W$ n'étant pas identifiable, la contagion se lit en bande
(chapitre~\ref{chap:partielle}). Le troisième, le plus souvent passé sous silence, est
l'incertitude de \textbf{modèle} : le choix de la \emph{famille} elle-même.

On la chiffre en recalculant le SCR sous des familles concurrentes. Sur l'axe de la
\textbf{dépendance}, à marges par pilier fixées, le SCR va de $5\,322$ M\euro{} (indépendance) à
$9\,806$ M\euro{} (comonotone, borne de Fréchet), soit un facteur $1{,}8$, la formule standard
Solvabilité~II et les copules gaussienne et de Gumbel se plaçant entre les deux ; notre cascade
dirigée ($8\,318$ M\euro) tombe dans cette bande, entre une dépendance gaussienne et une
dépendance de Gumbel. Sur l'axe de la \textbf{famille de queue}, l'écart est bien plus large :
d'une lognormale ($4\,330$ M\euro) à une GPD à queue lourde ($\xi=0{,}90$ ; $24\,016$ M\euro),
un facteur $5{,}5$. C'est le choix de la loi de queue qui domine, de loin, l'incertitude de
modèle --- et c'est précisément pourquoi le plafond de réassurance ($40$ M\euro) est un levier
distinct, qui ramène la famille lourde à $495$ M\euro : une réponse au risque de queue, non le
risque lui-même.

\begin{figure}[htbp]\centering
\includegraphics[width=\linewidth]{J4_bande_modele.png}
\caption{La bande de modèle, troisième étage d'incertitude. \textbf{(a)}~axe dépendance, à
marges fixées : la cascade se place entre les copules gaussienne et de Gumbel, dans une bande
$[5{,}3\,;9{,}8]$ Md\euro. \textbf{(b)}~axe famille de queue : la queue lourde domine
l'incertitude ($\times 5{,}5$). \textbf{(c)}~les trois étages (paramètre, identification,
modèle) se cumulent. Source : script~\texttt{48}.}
\label{fig:bande-modele}
\end{figure}

\begin{cle}
\textbf{Oui, chaque brique a des alternatives.} Fréquence (Poisson, Hawkes, Cox), sévérité
(lognormale, splicing, EVT bayésienne), dépendance (copules, réseaux, Hawkes multivarié),
mesure (TVaR, spectrale, robuste), agrégation (formule standard, Panjer) : on ne prétend pas
détenir \emph{la} méthode. La discipline est ailleurs : on a testé les concurrents les plus
sérieux, on \textbf{borne} l'effet du choix, et surtout le niveau bouge dans la bande mais
l'\textbf{ordre d'amorce des piliers ne bouge pas} ($P1>P4>P2>P3>P5$, invariant à la famille).
Le niveau est illustratif, l'ordre résiste : c'est, chiffrée sur un troisième axe, la thèse du
mémoire, et la réponse à qui demanderait \og et si vous aviez pris une autre méthode~?\fg.
\end{cle}

\section{De l'incertitude au chiffre reporté : VaR prédictive et VaR robuste}
\label{sec:predictive-robuste}

Quantifier l'incertitude (les trois étages) ne dit pas encore \emph{quel} nombre reporter.
Prendre le quantile $99{,}5\,\%$ d'une loi estimée en un point (\emph{plug-in}) est fragile,
car les données sont incertaines : \og{}le quantile d'un objet incertain est mauvais\fg{}
(C.~Hillairet). Deux réponses complémentaires, selon la posture (figures~\ref{fig:predictive}
et~\ref{fig:robuste}).

\textbf{La VaR prédictive} intègre l'incertitude de paramètre \emph{avant} de prendre le
quantile : on prend le quantile de la loi \emph{prédictive}, mélange des GPD sur la distribution
bootstrap de $(\xi,\sigma)$. Le résultat est instructif : à $99{,}5\,\%$ elle est \emph{quasi
égale} au \emph{plug-in} ($657$ contre $659$~M\euro) --- le risque d'estimation n'est pas un
biais de point à ce niveau, il ne le devient qu'en profondeur ($+6\,\%$ à $99{,}9\,\%$). La
fragilité est dans la \textbf{largeur} de la bande, pas dans le point.

\textbf{La VaR robuste} prend au contraire la \textbf{borne haute} de la VaR sur l'ensemble
d'ambiguïté (paramètre, puis modèle) : la VaR robuste à $95\,\%$ vaut $1\,026$~M\euro{}
($1{,}6\times$ le \emph{plug-in}), le pire-cas de famille lourde $1\,311$~M\euro. C'est la
posture prudente, où le niveau de robustesse est un \textbf{choix prudentiel explicite}.

\begin{figure}[htbp]\centering
\includegraphics[width=\linewidth]{J2_var_predictive.png}
\caption{VaR prédictive. \textbf{(a)}~le paramètre $\xi$ est incertain. \textbf{(b)}~les queues
\emph{plug-in} et prédictive ne divergent que loin dans la queue. \textbf{(c)}~le capital se lit
en profil et en bande. Source : script~\texttt{46}.}
\label{fig:predictive}
\end{figure}

\begin{figure}[htbp]\centering
\includegraphics[width=\linewidth]{J7_scr_robuste.png}
\caption{VaR robuste. \textbf{(a)}~la VaR est elle-même incertaine ; le robuste en prend la
borne haute. \textbf{(b)}~l'échelle des postures, du point ($659$) au pire-cas de modèle
($1\,311$~M\euro). Source : script~\texttt{51}.}
\label{fig:robuste}
\end{figure}

\begin{cle}
\textbf{Deux points honnêtes plutôt qu'un point trompeur.} La VaR prédictive est le point
\emph{honnête} (il embarque l'incertitude d'estimation) ; la VaR robuste est le point
\emph{prudent} (la borne haute sur l'ambiguïté). Le choix entre les deux relève de la politique
prudentielle, non du calcul. L'essentiel est de ne jamais reporter le seul \emph{plug-in} : le
SCR se cite avec sa bande, ou par un point qui assume l'incertitude, jamais nu.
\end{cle}

\begin{cle}
\textbf{La double lecture.} \og Insensible aux choix innocents, sensible aux leviers réels \fg :
le modèle n'est ni sur-ajusté ni fragile, il est \textbf{piloté par ce qui doit le piloter}.
Et la seule grande incertitude, la queue, est nommée et bornée, pas cachée. C'est ce qui
transforme \og je n'ai pas tout calibré \fg{} en \og je sais précisément de quoi mon SCR
dépend, et de quoi il ne dépend pas \fg.
\end{cle}

% =====================================================================
